% Start
\documentclass{beamer}
\usetheme{tudcd}
\iftutex
  \usepackage{fontspec}
\else
  \usepackage[T1]{fontenc}
  \usepackage[ngerman=ngerman-x-latest]{hyphsubst}
\fi
\usepackage[american]{babel}

\usepackage[backend=biber,sorting=none]{biblatex}
\addbibresource{citations.bib}

\usepackage{adjustbox}
\usepackage{pifont}% \ding

\usepackage{pdfpc}
\renewcommand*{\note}[1]{%
  \makebox[0pt]{% set width to 0pt
    \raisebox{0pt}[0pt][0pt]{% set height and depth to 0pt
      \pdfpcnote{#1}%
    }%
  }%
}
\def\lf{^^J}

\title[Crash-Tolerant Proxy]{A Crash-Tolerant Proxy Service for High-Availability Environments}
\author{\texorpdfstring{\texttt{\$AUTHOR}}{\$AUTHOR}}
\context{Bachelor Thesis}
\location{Dresden}
\date{2026-10-26}

\input{tikz/include/head.tex}
\urlstyle{same}
\begin{document}

\TUDCDsetup{
  color = {
    highlight-color=Weiss,
    secondary-color=Brilliantblau
  },
  logo = {
    color = white
  }
}
\maketitle
\TUDCDsetup{
  color = {
    highlight-color=Brilliantblau,
    secondary-color=Weiss
  },
  logo = {
    color = blue
  }
}

\begin{frame}{Motivation}
  \begin{columns}[T]
    \begin{column}{.5\textwidth}
      \begin{adjustbox}{max width=\textwidth}%
        \input{tikz/proxy-network.tex}%
      \end{adjustbox}%
    \end{column}%
    \begin{column}{0.4\textwidth}
      \begin{itemize}
        \item logging
        \item content-filtering
        \item authorization
        \item load-balancing
        \item network separation
        \item \ldots
      \end{itemize}%
    \end{column}
  \end{columns}%
  \note{%
    %* concurrent/asyncio\lf
    * offload tasks from client/server to proxy\lf
    * DBs, weil Idee von Prof. Lehner\lf
    * ..., TLS, abuse, re-routing, transform%
  }%
\end{frame}

\begin{frame}{Critical role in the network}
  \centering
  \begin{adjustbox}{max width=\textwidth}%
    \newif\ifcrash
    \crashtrue
    \input{tikz/proxy-network.tex}%
  \end{adjustbox}%

  \begin{minipage}{.6\textwidth}%
    \begin{alertblock}{Problem}%
      If the proxy terminates all its connections terminate.%
    \end{alertblock}%
  \end{minipage}%
  \note{%
    * planned restart/upgrade\lf
    * crashes/faults\lf
    * impact on adjacent services%
  }%
\end{frame}

\begin{frame}{Existing solutions}
  \begin{columns}[T]%
    \begin{column}{.55\textwidth}
      \begin{itemize}
      \item<1->
        protocol-level retries
        \only<1>{%
          \note{%
            * TCP-based\lf
            * TCP ist layer 4%
          }%
        }%

      \item<2->
        CRIU -- Checkpoint/Restore In Userspace \autocite{criu}

        \begin{itemize}
          \item
            checkpoint process trees
          \item
            supports TCP sockets with \texttt{TCP\_REPAIR}%
            \only<2>{%
              \note{%
                * TCP_REPAIR: read/write kernel buffers\lf
                * periodic vs reactive checkpoint%
              }%
            }%
        \end{itemize}%

      \item<2->
        VM checkpointing \autocite{clark2005live}

      \item<5->
        TCP-based replication (e.\,g., FT-TCP \autocite{zagorodnov2003engineering},
        ST-TCP \autocite{marwah2003tcp}, or CoRAL \autocite{aghdaie2005transparent})

        \begin{itemize}
          \item save incoming raw TCP traffic
          \item replay TCP traffic to replicate
        \end{itemize}%

\only<5>{%
        \note{%
          * TCP replay\lf
          * Logger wird point of failure\lf
          * CoRAL's raw sockets\lf
          * (Server) Fault Tolerant TCP\lf
          * Fetzer%
        }%
}%
      \end{itemize}
    \end{column}%
\onslide<5->{%
    \begin{column}{.4\textwidth}
      \vspace*{4em}%
      \flushright
      \begin{adjustbox}{max width=\textwidth}
        \shortstack{%
          \includegraphics[
            page=4,
            trim={20mm 55mm 116.9mm 180.4mm},
            clip,
          ]{st-tcp.pdf}\\
          \includegraphics[
            page=4,
            trim={20mm 39mm 116.9mm 233.2mm},
            clip,
          ]{st-tcp.pdf}%
        }%
      \end{adjustbox}\par
      \vspace{-.3em}%
      \scriptsize Credits: \textcite{marwah2003tcp}%
    \end{column}%
}%
  \end{columns}%

\only<3-4>{%
  \begin{tikzpicture}[remember picture, overlay]
    \node (foo) at (.5\textwidth,.5\textheight) {%
      \begin{minipage}{.8\textwidth}%
        \begin{block}{Limitations of checkpointing}%
          \begin{columns}[T]
            \begin{column}{0.45\textwidth}
              \begin{adjustbox}{max height=.7\textheight}%
                \begin{tikzpicture}[x=10em,y=-2em,>=Stealth]
                  \node[draw,anchor=south] at (0,-1) {\strut Host};
                  \node[draw,anchor=south] at (1,-1) {\strut Peer};
                  \draw (0,-1) -- (0,8.5);
                  \draw (1,-1) -- (1,8.5);

                  \draw[->] (1,0) -- node[above=.5ex,anchor=base] {\texttt{data[0:512]}} (0,0);
                  \draw[->] (0,0.5) to[bend left=-90,looseness=2.5] node[left] (transform) {%
                    \shortstack[r]{%
                      update kernel\\
                      buffers%
                    }%
                  } (0,1.5);
                  \draw[->] (0,2) -- node[above=.5ex,anchor=base] {\texttt{ACK 512}} (1,2);
                  \node[anchor=west] at (1,2) {seq\# 512};

                  \node[anchor=east] at (-2em,2.5) {\color{Gruen1}\shortstack[r]{checkpoint\\(seq\# 512)}};
                  \node[fill=Grau10] at (0,2.5) {};
                  \draw[dashed,Gruen1] (-2em,2.5) -- (2em,2.5);

                  \draw[->] (1,3.5) -- node[above=.5ex,anchor=base] {\texttt{data[512:1024]}} (0,3.5);
                  \draw[->] (0,4) to[bend left=-90,looseness=2.5] node[left] (transform) {%
                    \shortstack[r]{%
                      update kernel\\
                      buffers%
                    }%
                  } (0,5);
                  \draw[->] (0,5.5) -- node[above=.5ex,anchor=base] {\texttt{ACK 1024}} (1,5.5);
                  \node[anchor=west] at (1,5.5) {seq\# 1024};

                  \node[anchor=east] at (-2em,6) {\color{Rot1}{\Lightning} crash {\Lightning}};
                  \node[fill=Grau10] at (0,6) {};
                  \draw[dashed,Rot1] (-2em,6) -- (2em,6);
                \end{tikzpicture}
              \end{adjustbox}%
            \end{column}%
  \onslide<4>{%
            \begin{column}{0.45\textwidth}
              \begin{adjustbox}{max height=.7\textheight}%
                \begin{tikzpicture}[x=10em,y=-2em,>=Stealth]
                  \node[draw,anchor=south] at (0,-1) {\strut Host};
                  \node[draw,anchor=south] at (1,-1) {\strut Peer};
                  \draw (0,-1) -- (0,8.5);
                  \draw (1,-1) -- (1,8.5);

                  \node[anchor=east] at (-2em,2.5) {\color{Gruen1}\shortstack[r]{restore\\(seq\# 512)}};
                  \node[fill=Grau10] at (0,2.5) {};
                  \draw[dashed,Gruen1] (-2em,2.5) -- (2em,2.5);

                  \draw[->] (1,7) -- node[above=.5ex,anchor=base] {\texttt{data[1024:1536]}} (0,7);
                  \draw[->,Rot1] (0,8) -- node[above=.5ex,anchor=base] {\color{Rot1}{\Lightning} \texttt{ACK 512} {\Lightning}} (1,8);
                \end{tikzpicture}
              \end{adjustbox}%
            \end{column}%
  }%
          \end{columns}%
        \end{block}%
      \end{minipage}%
    };
    \fill[white,opacity=0.8] ({0,0} |- {foo.north}) rectangle (foo.south west);
  \end{tikzpicture}%
  \note{%
    * seq\# of peer\lf
    * no longer in send buffer%
  }%
}%
\end{frame}

\begin{frame}{Goal and error model}
  \begin{columns}%
    \begin{column}{.45\textwidth}%
      \begin{exampleblock}{\strut in-scope}%
        \begin{minipage}[c][.55\textheight]{\linewidth}%
          \begin{itemize}%
            \item
              transparent to clients and servers
            \item
              tolerate process termination

              \begin{itemize}
                \item implementation bugs
                \item by other processes
                \item by operating system (OOM killer)
              \end{itemize}
            \item
              isolation of components
          \end{itemize}%
        \end{minipage}%
      \end{exampleblock}%
    \end{column}%
    \begin{column}{.45\textwidth}
      \begin{alertblock}{\strut out-of-scope}
        \begin{minipage}[c][.55\textheight]{\linewidth}%
          \begin{itemize}%
            \item network interruptions
            \item misbehaving operating system
            \item hardware failures
            \item power loss
          \end{itemize}%
        \end{minipage}%
      \end{alertblock}%
    \end{column}%
  \end{columns}%
  \note{%
    * deterministic bugs = bad\lf
    * kill by signals\lf
    * isolate complex transformation\lf
    * unprivileged\lf
    * reusable%
  }%
\end{frame}

\begin{frame}{Connection lifecycle}
  \begin{columns}
    \begin{column}{.55\textwidth}
      \begin{itemize}
        \item
          transformation of traffic
        \item
          many connections in parallel

          \begin{itemize}
            \item asynchronous/non-blocking I/O
            \item interwoven control flow, i.e., coroutines
        \end{itemize}%
      \end{itemize}%
    \end{column}%
    \begin{column}{.4\textwidth}
      \begin{adjustbox}{max height=.8\textheight}%
        \input{tikz/connection-lifecycle-simple.tex}%
      \end{adjustbox}%
    \end{column}%
  \end{columns}%
  \note{%
    * 3 stages\lf
    * 2nd stage simplified\lf
    * alles = xfrm\lf
    * error page = xfrm\lf
    * not thread/process per connection\lf
    * 1st stage => fd copy%
  }%
\end{frame}

\begin{frame}{Redundant processes}
  \centering
  \begin{adjustbox}{max height=.8\textheight}%
    \begin{tikzpicture}[x=12em,y=-6em,>=Stealth]
      \node[draw,fill=gray!30,anchor=south] (fdtab1) at (0,0) {%
        \shortstack[l]{%
          %\textbullet\ IPC sockets\\
          \textbullet\ \textbf{listening sockets}\\
          \textbullet\ \textbf{shared memory}\\
          \textbullet\ \textbf{connection 1}\\
          \hphantom{\textbullet\ }\vphantom{connection 1}\ldots\\
          \textbullet\ \textbf{connection \textit{N}}%
        }%
      };
      \node[inner sep=0,anchor=north east] (label1) at (fdtab1.north west) {\textbf{\rlap{primary}\hphantom{\shortstack{primary\\backup}\hspace{.5em}}}};
      \node[draw,rounded corners,fit=(label1)(fdtab1),inner sep=.5em] (process1) {};

      \node[draw,fill=gray!30,anchor=north] (fdtab2) at (0,1) {%
        \shortstack[l]{%
          %\textbullet\ IPC sockets\\
          \textbullet\ \textbf{listening sockets}\\
          \textbullet\ \textbf{shared memory}\\
          \textbullet\ \textbf{connection 1}\\
          \hphantom{\textbullet\ }\vphantom{connection 1}\ldots\\
          \textbullet\ \textbf{connection \textit{N}}%
        }%
      };
      \node[inner sep=0,anchor=north east] (label2) at (fdtab2.north west) {\textbf{\rlap{backup}\hphantom{\shortstack{primary\\backup}\hspace{.5em}}}};
      \node[draw,rounded corners,fit=(label2)(fdtab2),inner sep=.5em] (process2) {};

      \draw[->] ({fdtab1} |- {process1.south})
        -- node[right] {\shortstack[l]{copy file\\descriptors}}
        ({fdtab2} |- {process2.north});

      \node[
        inner sep=.5em,
        path picture={
          \draw
            ($(path picture bounding box.north east) + (-.5\pgflinewidth,-.5\pgflinewidth)$)
            ++(-1em,0)
            -- ++(0,-1em)
            -- ++(1em,0)
            -- ++(-1em,1em)
            -- ($(path picture bounding box.north west) + (.5\pgflinewidth,-.5\pgflinewidth)$)
            -- ($(path picture bounding box.south west) + (.5\pgflinewidth,.5\pgflinewidth)$)
            -- ($(path picture bounding box.south east) + (-.5\pgflinewidth,.5\pgflinewidth)$)
            -- ($(path picture bounding box.north east) + (-.5\pgflinewidth,-.5\pgflinewidth) + (0,-1em)$);
        },
      ] (shmem) at (1,0.5) {\shortstack{\strut\\shared\\memory\\\strut}};

      \draw[->] (process1) -- (shmem);
      \draw[->] (process2) -- (shmem);
    \end{tikzpicture}%
  \end{adjustbox}%
  \note{%
    * kernel resources tied to fd\lf
    * Unix socket, SCM_RIGHTS\lf
    * shared memory%
  }%
\end{frame}

\begin{frame}{Redundant processes}
  \framesubtitle{%
      \vphantom{Shared file descriptor table}%
      \only<2->{Shared file descriptor table}%
  }

  \begin{columns}[T]
    \begin{column}{0.4\textwidth}
      \begin{alertblock}{Problem}
        \ttfamily
        \footnotesize
        \textcolor{Blau1}{int} fd = accept(listen\_fd,\newline
        \hphantom{\scalebox{16}{ }}NULL, NULL);\newline
        \textcolor{Gruen1}{// {\Lightning} CRASH {\Lightning}}\newline
        copy\_fd(fd);%
      \end{alertblock}%
    \end{column}%
    \begin{column}<2->{0.55\textwidth}%
      \begin{exampleblock}{Solution}%
        \centering
        \vspace{.5em}%
        \begin{adjustbox}{max width=\textwidth,max height=.8\textheight}%
          \input{tikz/clone-files.tex}%
        \end{adjustbox}%
        \vspace{.5em}%
      \end{exampleblock}%
    \end{column}%
  \end{columns}
\only<1>{%
  \note{%
    * auch für connect\lf
    * neue Verbindungen arguably unwichtig%
  }%
}%
\only<2->{%
  \note{%
    * CLONE_FILES Horsts Idee\lf
    * lost file descriptors%
  }%
}%
\end{frame}

\begin{frame}{Final process tree}
  \centering
  \begin{adjustbox}{max height=.8\textheight}%
    \newenvironment{secondslide}{\begin{uncoverenv}<2->}{\end{uncoverenv}}%
    \input{tikz/process-tree.tex}%
  \end{adjustbox}%
\only<1>{%
  \note{%
    * active/passive listener%
  }%
}%
\only<2->{%
  \note{%
    * worker scope\lf
    * privileges\lf
    * comparison nginx\lf
    * separate executable (execve vs clone)\lf
    * multiple workers\lf
    * interrupted shmem => Transactions%
  }%
}%
\end{frame}

\begin{frame}{Transactions}
  \begin{columns}
    \begin{column}{0.45\textwidth}
      \centering
      \begin{adjustbox}{max height=.8\textheight}%
        \begin{tikzpicture}[x=4em,y=-3em,>=Stealth]
          \node[circle,fill,inner sep=.2em] (start) at (0,0) {};
          \node[circle,draw,inner sep=.5em] (a) at (1,1) {};
          \node[circle,draw,inner sep=.5em] (b) at (1,2) {};
          \node[circle,draw,inner sep=.5em] (c) at (1,3) {};
          \node[rotate=90] (d) at (1,4) {\ldots};

          \draw[->] (start) -- (0,0.5) -- (a);
          \draw[->] (a) -- node[right] {\only<2->{\texttt{state := \ldots}}\only<1>{\hphantom{\texttt{state := \ldots}}}} (b);
          \draw[->] (b) -- node[right] {\only<2->{\texttt{state := \ldots}}} (c);
          \draw[->] (c) -- node[right] {\only<2->{\texttt{state := \ldots}}} (d);

\visible<3->{
          \node[anchor=east] at (0,0.5) {\texttt{state?}};
          \node[rotate=90] (e) at (0,3.5) {\ldots};
          \draw[dashed,->] (0,0.5) -- (e);
          \draw[dashed,->] (0,1.5) -- (b);
          \draw[dashed,->] (0,2.5) -- (c);
}
        \end{tikzpicture}%
      \end{adjustbox}%
    \end{column}%
    \begin{column}{0.5\textwidth}
      \begin{itemize}
        \item split into idempotent steps
        \item<2-> update \texttt{state} atomically
        \item<3-> continue with interrupted step
      \end{itemize}%
    \end{column}%
  \end{columns}%
  \note{%
    * not all of ACID, no isolation, durability\lf
    * explain idempotent\lf
    * Vgl. basis block\lf
    * atomic = single instruction\lf
    * stdatomic.h, C11%
  }%
\end{frame}

\begin{frame}{Transactions}{I/O}
  \centering
  \begin{adjustbox}{max height=.8\textheight}%
    \begin{tikzpicture}[x=10em,y=-2em,>=Stealth]
      \node[draw,anchor=south] at (0,-1) {\strut kernel};
      \node[draw,anchor=south] at (1,-1) {\strut userspace};
      \draw (0,-1) -- (0,8.5);
      \draw (1,-1) -- (1,8.5);

      \draw[->] (-0.5,0) -- node[above=.5ex,anchor=base] {\texttt{data}} (0,0);
      \draw[->] (0,0.5) to[bend left=90,looseness=2.5] node[right] (transform) {%
        \shortstack[l]{%
          update kernel\\
          buffers%
        }%
      } (0,1.5);
      \draw[->] (0,2) -- node[above=.5ex,anchor=base] {\texttt{ACK}} (-0.5,2);

      \draw[->] (1,3.5) -- node[above=.5ex,anchor=base] {\manpage{2}{recv}} (0,3.5);
      \draw[->] (0,4) to[bend left=90,looseness=2.5] node[right] (transform) {%
        \shortstack[l]{%
          update kernel\\
          buffers%
        }%
      } (0,5);
      \draw[->] (0,6) -- node[above=.5ex,anchor=base] {\texttt{return}} (1,6);

\visible<2->{
      \node[fill=white] at (0,6.5) {};
      \node[fill=white] at (1,6.5) {};
      \draw[dashed,Rot1] (1,6.5) ++(5em,0) -- node[fill=white] (sigkill) {\color{Rot1}\texttt{SIGKILL}} (-5em,6.5);
}

      \draw[->] (1,7) to[bend left=90,looseness=2.5] node[right] (transform) {%
        \shortstack[l]{%
          update userspace\\
          buffers%
        }%
      } (1,8);
      \draw[opacity=0] (0,7.5) -- ++($(1,4.5) - (transform.east)$);
    \end{tikzpicture}
  \end{adjustbox}%
  \note{%
    * kernel buffers bound to file descriptor\lf
    * buffer in shared memory\lf
    * syscall signal-atomic\lf
    * Vgl. send%
  }%
\only<2>{%
  \note{%
    * lost return value\lf
    * calling conventions\lf
    * => recvmmsg%
  }%
}%
\end{frame}

\begin{frame}{Transactions}{\manpage{2}{recvmmsg}, \manpage{2}{sendmmsg}}
  \begin{columns}
    \begin{column}{0.55\textwidth}
      \ttfamily
      \footnotesize
      \textcolor{Blau1}{struct} mmsghdr \{\newline
      \hphantom{ \ \ \ }\textcolor{Blau1}{struct} msghdr msg\_hdr;\newline
      \hphantom{ \ \ \ }\textcolor{Blau1}{unsigned int} msg\_len;\newline
      \};\newline
      \newline
      \textcolor{Blau1}{int} recvmmsg(\textcolor{Blau1}{int} sockfd,\newline
      \hphantom{ \ \ \ \ \ \ \ \ \ \ \ \ }\textcolor{Blau1}{struct} mmsghdr *msgvec,\newline
      \hphantom{ \ \ \ \ \ \ \ \ \ \ \ \ }\textcolor{Blau1}{unsigned int} n,\newline
      \hphantom{ \ \ \ \ \ \ \ \ \ \ \ \ }\textcolor{Blau1}{int} flags,\newline
      \hphantom{ \ \ \ \ \ \ \ \ \ \ \ \ }\textcolor{Blau1}{struct} timespec *timeout);%
    \end{column}%
    \begin{column}{0.35\textwidth}%
      \begin{adjustbox}{max height=.8\textheight}%
        \begin{tikzpicture}[x=10em,y=-2em,>=Stealth]
          \node[draw,anchor=south] at (0,-1) {\strut kernel};
          \node[draw,anchor=south] at (1,-1) {\strut userspace};
          \draw (0,-1) -- (0,8);
          \draw (1,-1) -- (1,8);

          \draw[->] (-0.5,0) -- node[above=.5ex,anchor=base] {\texttt{data}} (0,0);
          \draw[->] (0,0.5) to[bend left=90,looseness=2.5] node[right] (transform) {%
            \shortstack[l]{%
              update kernel\\
              buffers%
            }%
          } (0,1.5);
          \draw[->] (0,2) -- node[above=.5ex,anchor=base] {\texttt{ACK}} (-0.5,2);

          \draw[->] (1,3.5) -- node[above=.5ex,anchor=base] {\manpage{2}{recvmmsg}} (0,3.5);
          \draw[->] (0,4) to[bend left=90,looseness=2.5] node[right] (transform) {%
            \shortstack[l]{%
              update kernel\\
              buffers%
            }%
          } (0,5);
\visible<2->{
          \draw[->] (0,5.5) to[bend left=90,looseness=2.5] node[right] (transform) {%
            \shortstack[l]{%
              update\\
              \texttt{\textcolor{Blau1}{struct} mmsghdr}%
            }%
          } (0,6.5);
          \draw[->] (0,7.5) -- node[above=.5ex,anchor=base] {\texttt{return}} (1,7.5);
}
        \end{tikzpicture}%
      \end{adjustbox}%
    \end{column}%
  \end{columns}%
  \note{%
    * multiple recvmsg\lf
    * retval in msg_len\lf
    * struct mmsghdr in shared memory\lf
    * syscall is signal-atomic\lf
    * possibly lost -1%
  }%
\end{frame}

\begin{frame}{Transactions}{\manpage{2}{recvmmsg}, \manpage{2}{sendmmsg}}
  \centering
  \resizebox{!}{.8\textheight}{%
    \begin{tikzpicture}[x=6em,y=-2em,>=Stealth]
      \node[circle,fill,inner sep=.2em] (start) at (0,-0.5) {};

      \node[rounded corners,draw,inner sep=.5em,anchor=north] (a) at (1,1) {\shortstack{write sentinel\\value (\texttt{-1}) to\\\texttt{\textcolor{Blau1}{struct} mmsghdr}}};
      \draw[->] (start) -- ($(a.north) + (-1,-0.5)$) -- (a);

\visible<2->{
      \node[rounded corners,draw,inner sep=.5em,anchor=north] (b) at ($(a.south) + (0,1)$) {\itshape \shortstack{sentinel value\\was overwritten}};
      \draw[->] (a) -- node[right] {\texttt{recvmmsg()}} (b);
      \draw[dashed,->] (0,0.5) -- ($(b.north) + (-1,-0.5)$) -- (b);
}

\visible<3->{
      \node[rounded corners,draw,inner sep=.5em,anchor=north] (c) at ($(b.south) + (0,1)$) {\shortstack{update buffers from\\\texttt{\textcolor{Blau1}{struct} mmsghdr}}};
      \draw[->] (b) -- node[right] {\texttt{state := APPEND}} (c);
      \draw[dashed,->] ($(b.north) + (-1,-0.5)$) -- ($(c.north) + (-1,-0.5)$) -- (c);

      \node[circle,draw,inner sep=.5em,anchor=north] (d) at ($(c.south) + (0,1)$) {};
      \node[circle,draw,inner sep=.35em] at (d) {};
      \draw[->] (c) -- node[right] {\texttt{state := DONE}} (d);
}
    \end{tikzpicture}%
  }%
  \note{%
    * sentinel\lf
    * repeated syscall = bad\lf
    * signal-atomic%
  }%
\end{frame}

\begin{frame}{Transactions}{Transformation}
  \begin{columns}
    \begin{column}{0.6\textwidth}
      \begin{adjustbox}{max width=\textwidth,max height=.8\textheight}%
        \input{tikz/transform.tex}%
      \end{adjustbox}%
    \end{column}%
    \begin{column}{0.3\textwidth}
\visible<2->{%
      \centering
      \resizebox{!}{.8\textheight}{%
        \begin{tikzpicture}[x=6em,y=-2em,>=Stealth]
          \node[circle,fill,inner sep=.2em] (start) at (0,-0.5) {};

          \node[rounded corners,draw,fill=white,inner sep=.5em,anchor=north] (a) at (1,1) {\shortstack{copy all\\five buffers}};
          \draw[->] (start) -- ($(a.north) + (-1,-0.5)$) -- (a);

          \node[rounded corners,draw,fill=white,inner sep=.5em,anchor=north] (b) at ($(a.south) + (0,0.5)$) {\shortstack{\texttt{transform()}\\on copies}};
          \draw[->] (a) -- (b);

          \node[rounded corners,draw,fill=white,inner sep=.5em,anchor=north] (c) at ($(b.south) + (0,1)$) {\shortstack{swap buffers\\and copies}};
          \draw[->] (b) -- node[right] {\texttt{state := SWAP}} (c);

          \draw[dashed,->] (0,0.5) -- ($(c.north) + (-1,-0.5)$) -- (c);

          \node[circle,draw,fill=white,inner sep=.5em,anchor=north] (d) at ($(c.south) + (0,1)$) {};
          \node[circle,draw,inner sep=.35em] at (d) {};
          \draw[->] (c) -- node[right] {\texttt{state := DONE}} (d);
        \end{tikzpicture}%
      }%
}%
    \end{column}%
  \end{columns}%
  \note{%
    * all or nothing\lf
    * double buffering\lf
    * ring buffers not duplicate\lf
    * (skip transaction)%
  }%
\end{frame}

\begin{frame}<1>[label=design-summary]
  \frametitle{%
    \only<2>{Recap: }Design Summary%
  }%
  \begin{columns}%
    \begin{column}{0.5\textwidth}%
      \begin{adjustbox}{max height=.7\textheight}%
        \newif\ifchecked
        \checkedtrue
        \input{tikz/process-tree.tex}%
      \end{adjustbox}%
    \end{column}%
    \begin{column}{0.4\textwidth}%
      \hfill
      \begin{adjustbox}{max height=.7\textheight}%
        \begin{tikzpicture}[x=8em,y=-2em,>=Stealth]
          \node[anchor=south,draw] at (0,-1) {\strut Client};
          \node[anchor=south,draw] at (1,-1) {\strut Proxy};
          \node[anchor=south,draw] at (2,-1) {\strut Server};

          \fill[Blau2] (-0.25,-0.75) rectangle (2.25,1.25);
          \fill[Gruen2] (-0.25,1.75) rectangle (2.25,6.25);

          \draw (0,-1) -- (0,8.5);
          \draw (1,-1) -- (1,8.5);
          \draw (2,-1) -- (2,8.5);

          \draw[->] (0,0) -- node[above=.5ex,anchor=base] {\manpage{2}{accept}} (1,0);
          \draw[->] (1,0.5) -- node[above=.5ex,anchor=base] {\manpage{2}{connect}} (2,0.5);

          \node[fill=white] at (0,1.5) {};
          \node[fill=white] at (1,1.5) {};
          \node[fill=white] at (2,1.5) {};
          \draw[dashed] (-0.25,1.5) -- (2.25,1.5);

          \draw[->] (0,2.5) -- node[above=.5ex,anchor=base] {\manpage{2}{recv}} (1,2.5);
          \draw[->] (2,3) -- node[above=.5ex,anchor=base] {\manpage{2}{recv}} (1,3);
          \draw[->] (1,3.5) to[bend left=90,looseness=3] node[right] (xfrm) {transform} (1,4.5);
          \draw[->] (1,5) -- node[above=.5ex,anchor=base] {\manpage{2}{send}} (0,5);
          \draw[->] (1,5.5) -- node[above=.5ex,anchor=base] {\manpage{2}{send}} (2,5.5);

          \node[fill=white] at (0,6.5) {};
          \node[fill=white] at (1,6.5) {};
          \node[fill=white] at (2,6.5) {};
          \draw[dashed] (-0.25,6.5) -- (2.25,6.5);

          \draw[->] (1,7.5) -- node[above=.5ex,anchor=base] {\manpage{2}{close}} (2,7.5);
          \draw[->] (1,8) -- node[above=.5ex,anchor=base] {\manpage{2}{close}} (0,8);
        \end{tikzpicture}%
      \end{adjustbox}%
    \end{column}%
  \end{columns}%
  \note{%
    * fd handling in listener\lf
    * blau = fd\lf
    * grün = transactions\lf
    * close in beiden\lf
    * everyone holds ref%
  }%
\end{frame}

\begin{frame}{Evaluation}{Error injection}
  \begin{columns}%
    \begin{column}{0.5\textwidth}%
      \ttfamily
      \footnotesize
      \textcolor{Blau1}{int} conn\_fd = accept(...);\newline
      \newline
      \textcolor{Rot1}{libcrash\_accept\_post}(conn\_fd, ...);\newline
      \newline
      \textcolor{Gruen1}{// Save file descriptor.}\newline
      conn->downstream.fd[0] = conn\_fd;%
    \end{column}%
\visible<2->{%
    \begin{column}{0.4\textwidth}%
      \begin{adjustbox}{max height=.8\textheight}%
        \begin{tikzpicture}[x=8em,y=-2em,>=Stealth]
          \node[anchor=south,draw] at (0,-1) {\strut Client};
          \node[anchor=south,draw] at (1,-1) {\strut Proxy};
          \node[anchor=south,draw] at (2,-1) {\strut Server};

          \draw (0,-1) -- (0,8.5);
          \draw (1,-1) -- (1,8.5);
          \draw (2,-1) -- (2,8.5);

          \draw[->,Rot1] (0,0) -- node[above=.5ex,anchor=base] {\manpage{2}{accept}} (1,0);
          \draw[->,Rot1] (1,0.5) -- node[above=.5ex,anchor=base] {\manpage{2}{connect}} (2,0.5);

          \node[fill=white] at (0,1.5) {};
          \node[fill=white] at (1,1.5) {};
          \node[fill=white] at (2,1.5) {};
          \draw[dashed] (-0.25,1.5) -- (2.25,1.5);

          \draw[->,Rot1] (0,2.5) -- node[above=.5ex,anchor=base] {\manpage{2}{recv}} (1,2.5);
          \draw[->,Rot1] (2,3) -- node[above=.5ex,anchor=base] {\manpage{2}{recv}} (1,3);
          \draw[->] (1,3.5) to[bend left=90,looseness=3] node[right] {transform} (1,4.5);
          \draw[->,Rot1] (1,5) -- node[above=.5ex,anchor=base] {\manpage{2}{send}} (0,5);
          \draw[->,Rot1] (1,5.5) -- node[above=.5ex,anchor=base] {\manpage{2}{send}} (2,5.5);

          \node[fill=white] at (0,6.5) {};
          \node[fill=white] at (1,6.5) {};
          \node[fill=white] at (2,6.5) {};
          \draw[dashed] (-0.25,6.5) -- (2.25,6.5);

          \draw[->] (1,7.5) -- node[above=.5ex,anchor=base] {\manpage{2}{close}} (2,7.5);
          \draw[->] (1,8) -- node[above=.5ex,anchor=base] {\manpage{2}{close}} (0,8);
        \end{tikzpicture}
      \end{adjustbox}%
    \end{column}%
}%
  \end{columns}%
  \note{%
    * als shared library\lf
    * 11 hooks\lf
    * unit tests (I/O)\lf
    * schwurbel xfrm%
  }%
\end{frame}

\begin{frame}{Evaluation}{Setup}
  \begin{center}%
    \centering
    \begin{adjustbox}{max height=.8\textheight}%
      \input{tikz/veth.tex}%
    \end{adjustbox}%
  \end{center}%

  \begin{itemize}
    \item Intel Core i9-13900K
    \item 128\,GiB RAM
    \item Linux kernel 6.12.101
  \end{itemize}
  \note{%
    * bisher proto-agnostic\lf
    * HTTP easy & ecosystem\lf
    * explain baseline architecture\lf
    * why HAProxy%
  }%
\end{frame}

\begin{frame}{Evaluation}{No-op transformation}
  \centering
  \begin{adjustbox}{max height=.75\textheight}%
    \newif\ifnolegend
    \nolegendtrue
    \input{tikz/vegeta-transform_nop-1M-100c.tex}%
  \end{adjustbox}

  1\,MiB response, 100 concurrent connections

  \textcolor {green!75!black}{\rule{1ex}{1ex}} nginx |
  \textcolor{orange!75!black}{\rule{1ex}{1ex}} HAProxy (TCP) |
  \textcolor  {blue!75!black}{\rule{1ex}{1ex}} baseline variant |
  \textcolor   {red!75!black}{\rule{1ex}{1ex}} crash-tolerant variant%
  \note{%
    * 1 MiB, 100 conn\lf
    * 600 seconds\lf
    * median: 2.84ms, 68.26ms, 32.62ms, 36.86ms\lf
    * dedicated listener\lf
    * steps?%
  }%
\end{frame}

\begin{frame}{Evaluation}{HTTP header transformation}
  \centering
  \begin{adjustbox}{max height=.75\textheight}%
    \newif\ifnolegend
    \nolegendtrue
    \input{tikz/vegeta-transform_headers-1M-10c.tex}%
  \end{adjustbox}

  1\,MiB response, 10 concurrent connections

  \textcolor{orange!75!black}{\rule{1ex}{1ex}} HAProxy (HTTP) |
  \textcolor  {blue!75!black}{\rule{1ex}{1ex}} baseline variant |
  \textcolor   {red!75!black}{\rule{1ex}{1ex}} crash-tolerant variant%
  \note{%
    * add/remove HTTP headers\lf
    * 1 MiB, 10 conn\lf
    * 600 seconds\lf
    * bad HTTP\lf
    * median: 3.9ms, 5.85ms, 5.89ms%
  }%
\end{frame}

\begin{frame}{Evaluation}{Recovery}
  \vfill
  \centering
  \begin{adjustbox}{max height=.8\textheight}%
    \input{tikz/fail-recv.tex}%
  \end{adjustbox}%
  \note{%
    * N-1 conns\lf
    * N-th fails recvmmsg\lf
    * dumb recovery (/proc/self/fd)\lf
    * all connections because only one worker\lf
    * bottleneck IPC%
  }%
\end{frame}

\begin{frame}{Conclusion}
  \begin{itemize}
  \item
    successful proof of concept
  \item
    runtime overhead: between 0.3\,\% and 21.8\,\%
  \item
    memory overhead: 424\,B per connection (+ passive listener)
  \item
    crash recovery covers live-upgrade
  \item<2->
    BUT corruption of shared memory possible
  \end{itemize}

  \note{%
    * depending on xfrm & size\lf
    * many transactions/connections > overhead\lf
    * upgrade%
  }%
\end{frame}

\begin{frame}{Future work}
  \begin{columns}%
    \begin{column}{0.3\textwidth}%
      \begin{itemize}
        \item protect shared memory
        \item multiple workers
        \item session state
        \item \manpage{7}{io\_uring}
        \pause
        \item manual \texttt{ACK}
      \end{itemize}%
    \end{column}%
    \begin{column}{0.6\textwidth}%
      \hfill
      \begin{adjustbox}{max height=.8\textheight}%
        \input{tikz/ack-recv.tex}%
      \end{adjustbox}%
    \end{column}%
  \end{columns}%
  \note{%
    * io_uring mit accept%
  }%
\end{frame}

\againframe<2>{design-summary}

\begin{frame}
  \frametitle{Bibliography}
  {%
    \renewcommand*{\bibfont}{\normalfont\scriptsize}%
    \printbibliography
  }%
\end{frame}

\end{document}
